% F082｜04 例题：从初中算例到高考速答
\clearpage
\begin{examplebox}
	\textbf{\textcolor{CExample}{例 0：先认清“相减”可以得到负数}}
\begin{center}\begin{tikzpicture}[font=\small,>=stealth,line join=round]
\draw[gray!40,rounded corners=3pt] (0,0) rectangle (14.2,7.8);\node at (7.1,7.2) {同一对固定点，不同动点可以给出负差与正差};
\begin{scope}[shift={(2.0,3)},x=0.85cm,y=0.85cm]\draw[->,gray] (-0.6,0)--(2.6,0) node[right]{$x$};\draw[->,gray] (0,-.35)--(0,2.8) node[left]{$y$};\draw[orange!85!black,line width=1.6pt] (0,2)--(0,0);\draw[violet,thick] (2,1)--(0,0);\draw[teal!80!black,thick] (0,2)--(2,1);\fill[teal!80!black] (0,2) circle(2.2pt) node[left=6pt]{$A$};\fill[teal!80!black] (2,1) circle(2.2pt) node[above right=4pt]{$B$};\fill[blue] (0,0) circle(2.2pt) node[below=4pt]{$P$};\end{scope}\node at (3.6,6.2) {$x=0$：差为负};\node at (3.6,1.5) {$2-\sqrt5<0$};\begin{scope}[shift={(8.0,3)},x=0.85cm,y=0.85cm]\draw[->,gray] (-0.6,0)--(4.6,0) node[right]{$x$};\draw[->,gray] (0,-.35)--(0,2.8) node[left]{$y$};\draw[orange!85!black,line width=1.6pt] (0,2)--(4,0);\draw[violet,thick] (2,1)--(4,0);\draw[teal!80!black,thick] (0,2)--(2,1);\fill[teal!80!black] (0,2) circle(2.2pt) node[left=6pt]{$A$};\fill[teal!80!black] (2,1) circle(2.2pt) node[above right=4pt]{$B$};\fill[blue] (4,0) circle(2.2pt) node[below=4pt]{$P$};\end{scope}\node at (9.6,6.2) {$x=4$：差为正};\node at (9.6,1.5) {$2\sqrt5-\sqrt5=\sqrt5$};
\node[text width=12.8cm,align=center] at (7.1,.65) {橙色为第一段 $PA$，紫色为第二段 $PB$；必须按这个顺序相减。};
\end{tikzpicture}\end{center}



	固定 $A(0,2)$、$B(2,1)$、$P(x,0)$，
	\[
		f(x)=\sqrt{x^2+4}-\sqrt{(x-2)^2+1}.
	\]
	当 $x=0$ 时，$PA=2$，$PB=\sqrt5$，所以
	\[
		f(0)=2-\sqrt5<0.
	\]
	当 $x=4$ 时，$PA=\sqrt{20}=2\sqrt5$，$PB=\sqrt5$，所以
	\[
		f(4)=\sqrt5>0.
	\]
	\textbf{同一对定点，因为动点的位置改变，距离差可能从负数变成正数。} 求值域时绝不能只求差的绝对值。
	
\end{examplebox}
\clearpage
\begin{examplebox}
	\textbf{\textcolor{CExample}{例 1：$a>b$，有最大值但没有最小值}}
\begin{center}\begin{tikzpicture}[font=\small,>=stealth,line join=round]
\draw[gray!40,rounded corners=3pt] (0,0) rectangle (14.2,10.2);\node at (7.1,9.6) {第一项高度更大：先增后减，只有最大值能取到};
\node[blue] at (7.1,8.65) {$\displaystyle f(x)=\sqrt{x^2+4}-\sqrt{(x-2)^2+1}$};\begin{scope}[shift={(3.1,4.6)},x=0.36cm,y=1cm]\draw[->,gray] (-6,0)--(14.5,0) node[right]{$x$};\draw[->,gray] (0,-2.8)--(0,3.4000000000000004) node[left]{$y$};\draw[gray,dashed] (-6,-2)--(14,-2);\node[left=3pt,fill=white,inner sep=1pt] at (0,-2) {$-2$};\draw[gray,dashed] (-6,2)--(14,2);\node[left=3pt,fill=white,inner sep=1pt] at (0,2) {$2$};\draw[blue,very thick,domain=-6:14,samples=240] plot(\x,{sqrt(\x*\x+4)-sqrt((\x-2)^2+1)});\draw[gray,dotted] (4,0)--(4,2.23606797749979);\fill[blue] (4,2.23606797749979) circle(2.3pt);\node[below=4pt] at (4,0){$4$};\node[blue,above=9pt] at (4,2.23606797749979) {$(4,\sqrt5)$};\end{scope}\node[anchor=west,align=left] at (9,4.5) {$a=2>b=1$\\[12pt]左端：趋近 $-2$\\[10pt]右端：趋近 $2$\\[12pt]最大值：$\sqrt5$\\[12pt]完整值域：$(-2,\sqrt5]$};
\node[text width=12.8cm,align=center] at (7.1,.65) {虚线是无限远处靠近的高度；不是有限定义域的端点。};
\end{tikzpicture}\end{center}



	求
	\[
		f(x)=\sqrt{x^2+4}-\sqrt{(x-2)^2+1},\quad x\in\mathbb R
	\]
	的值域。

	\textbf{第 1 步：} $m=0,n=2,a=2,b=1$，所以 $a>b$。
	\[
		d=|0-2|=2,\qquad
		D=\sqrt{(0-2)^2+(2-1)^2}=\sqrt5.
	\]
	\textbf{第 2 步：} 直接用 $a>b$ 的结论：
	\[
		\boxed{f(x)\in(-2,\sqrt5]}.
	\]
	\textbf{第 3 步：检验右端真的能取到。}
	\[
		x_0=\frac{1\cdot0-2\cdot2}{1-2}=4,
		\qquad f(4)=2\sqrt5-\sqrt5=\sqrt5.
	\]
	左端 $-2$ 只能在 $x$ 往负无穷远移动时越来越接近，\textbf{不能写成 $[-2,\sqrt5]$}。
\end{examplebox}
\clearpage
\begin{examplebox}
	\textbf{\textcolor{CExample}{例 2：$a<b$，有最小值但没有最大值}}
\begin{center}\begin{tikzpicture}[font=\small,>=stealth,line join=round]
\draw[gray!40,rounded corners=3pt] (0,0) rectangle (14.2,10.2);\node at (7.1,9.6) {第一项高度更小：先减后增，只有最小值能取到};
\node[blue] at (7.1,8.65) {$\displaystyle g(x)=\sqrt{x^2+1}-\sqrt{(x-2)^2+4}$};\begin{scope}[shift={(3.1,4.6)},x=0.36cm,y=1cm]\draw[->,gray] (-6,0)--(14.5,0) node[right]{$x$};\draw[->,gray] (0,-2.8)--(0,3.4000000000000004) node[left]{$y$};\draw[gray,dashed] (-6,-2)--(14,-2);\node[left=3pt,fill=white,inner sep=1pt] at (0,-2) {$-2$};\draw[gray,dashed] (-6,2)--(14,2);\node[left=3pt,fill=white,inner sep=1pt] at (0,2) {$2$};\draw[blue,very thick,domain=-6:14,samples=240] plot(\x,{sqrt(\x*\x+1)-sqrt((\x-2)^2+4)});\draw[gray,dotted] (-2,0)--(-2,-2.23606797749979);\fill[blue] (-2,-2.23606797749979) circle(2.3pt);\node[below=4pt] at (-2,0){$-2$};\node[blue,below=9pt] at (-2,-2.23606797749979) {$(-2,-\sqrt5)$};\end{scope}\node[anchor=west,align=left] at (9,4.5) {$a=1<b=2$\\[12pt]左端：趋近 $-2$\\[10pt]右端：趋近 $2$\\[12pt]最小值：$-\sqrt5$\\[12pt]完整值域：$[-\sqrt5,2)$};
\node[text width=12.8cm,align=center] at (7.1,.65) {橙紫两项顺序固定；这里的最低点为负，不能把距离差改成绝对值。};
\end{tikzpicture}\end{center}



	求
	\[
		g(x)=\sqrt{x^2+1}-\sqrt{(x-2)^2+4}
	\]
	的值域。

	这里 $a=1<b=2$，$m=0,n=2$，故
	\[
		d=2,\qquad D=\sqrt{2^2+(1-2)^2}=\sqrt5.
	\]
	因此\textbf{这次恰好和例 1 的开闭方向相反}：
	\[
		\boxed{g(x)\in[-\sqrt5,2)}.
	\]
	取最小值的位置为
	\[
		x_0=\frac{2\cdot0-1\cdot2}{2-1}=-2.
	\]
	验证：$g(-2)=\sqrt5-\sqrt{20}=-\sqrt5$。最大值不存在，$2$ 只能从下方越来越接近。

	\textbf{一眼速答：}高度较低的第一点 $A$ 让差达到\textbf{负的} $AB$；高度较高的第二点 $B$ 不允许正的 $AB$ 也达到。
\end{examplebox}
\clearpage
\begin{examplebox}
	\textbf{\textcolor{CExample}{例 3：$a=b$，两个端点都取不到}}
\begin{center}\begin{tikzpicture}[font=\small,>=stealth,line join=round]
\draw[gray!40,rounded corners=3pt] (0,0) rectangle (14.2,9.2);\node at (7.1,8.6) {高度相同：一直增大，两端都只能靠近};
\node[blue] at (7.1,7.65) {$\displaystyle h(x)=\sqrt{x^2+4}-\sqrt{(x-3)^2+4}$};\begin{scope}[shift={(4.1,4.0)},x=0.3130434782608696cm,y=0.75cm]\draw[->,gray] (-10,0)--(13.5,0) node[right]{$x$};\draw[->,gray] (0,-3.6)--(0,3.8000000000000003) node[left]{$y$};\draw[gray,dashed] (-10,-3)--(13,-3);\node[left=3pt,fill=white,inner sep=1pt] at (0,-3) {$-3$};\draw[gray,dashed] (-10,3)--(13,3);\node[left=3pt,fill=white,inner sep=1pt] at (0,3) {$3$};\draw[blue,very thick,domain=-10:13,samples=240] plot(\x,{sqrt(\x*\x+4)-sqrt((\x-3)^2+4)});\end{scope}\node[anchor=west,align=left] at (9,4.5) {$a=b=2,\ m=0,n=3$\\[12pt]对每个实数 $x$\\[10pt]都有 $-3<h(x)<3$\\[12pt]$h(3/2)=0$\\[12pt]完整值域：$(-3,3)$};
\node[text width=12.8cm,align=center] at (7.1,.65) {两条虚线均不属于函数图像；同高度且不同中心时没有最大、最小值。};
\end{tikzpicture}\end{center}



	求
	\[
		f(x)=\sqrt{x^2+4}-\sqrt{(x-3)^2+4}
	\]
	的值域。

	两定点 $A(0,2)$、$B(3,2)$ \textbf{高度一样}，$d=3$，所以
	\[
		\boxed{f(x)\in(-3,3)}.
	\]
	例如 $x=\tfrac32$ 时，两个根号完全相等，$f(\tfrac32)=0$。

	为什么不能取到 $\pm3$？两点连线是 $y=2$，动点始终位于 $y=0$；这两条水平线永远不会相交。只有把 $x$ 移得越来越远，距离差才会越来越接近 $\pm3$。
	
\end{examplebox}
\clearpage
\begin{examplebox}
	\textbf{\textcolor{CExample}{例 4：$m=n$，两个根号共用同一个中心}}
\begin{center}\begin{tikzpicture}[font=\small,>=stealth,line join=round]
\draw[gray!40,rounded corners=3pt] (0,0) rectangle (14.2,10.2);\node at (7.1,9.6) {例4：同一中心的距离差，远离中心时趋近零};
\node[blue] at (7.1,8.65) {$\displaystyle f(x)=\sqrt{(x-1)^2+9}-\sqrt{(x-1)^2+1}$};\begin{scope}[shift={(4,3.0)},x=0.5142857142857143cm,y=1.3cm]\draw[->,gray] (-6,0)--(8.5,0) node[right]{$x$};\draw[->,gray] (0,-0.6)--(0,3.4000000000000004) node[left]{$y$};\draw[gray,dashed] (-6,0)--(8,0);\node[left=3pt,fill=white,inner sep=1pt] at (0,0) {$0$};\draw[blue,very thick,domain=-6:8,samples=240] plot(\x,{sqrt((\x-1)^2+9)-sqrt((\x-1)^2+1)});\draw[gray,dotted] (1,0)--(1,2);\fill[blue] (1,2) circle(2.3pt);\node[below=4pt] at (1,0){$1$};\node[blue,above=9pt] at (1,2) {$(1,2)$};\end{scope}\node[anchor=west,align=left] at (9,4.5) {$m=n=1,\ a=3>b=1$\\[12pt]对每个实数 $x$：$f(x)>0$\\[12pt]值域 $(0,2]$\\[12pt]若两根号完全相同\\[9pt]则恒为 $0$，值域 $\{0\}$。};
\node[text width=12.8cm,align=center] at (7.1,.65) {蓝线表示题给函数；虚线 $y=0$ 是其渐近高度。相同根号的差才恒为零。};
\end{tikzpicture}\end{center}



	求
	\[
		f(x)=\sqrt{(x-1)^2+9}-\sqrt{(x-1)^2+1}
	\]
	的值域。

	因为 $m=n=1$，所以 $d=0$；$a=3>b=1$，故 $D=2$，直接得到
	\[
		\boxed{f(x)\in(0,2]}.
	\]
	在 $x=1$ 时，$f(1)=3-1=2$。当 $x$ 越来越远离 $1$ 时，两段距离差越来越小，却始终\textbf{严格大于 $0$}。

	如果把两个根号都改成同一个式子，例如
	\[
		\sqrt{(x-1)^2+9}-\sqrt{(x-1)^2+9},
	\]
	结果始终是 $0$，值域不是一个区间，而是 $\boxed{\{0\}}$。
\end{examplebox}
\clearpage
\begin{examplebox}
	\textbf{\textcolor{CExample}{例 5：隐藏在二次式中的距离差（配方题）}}
\begin{center}\begin{tikzpicture}[font=\small,>=stealth,line join=round]
\draw[gray!40,rounded corners=3pt] (0,0) rectangle (14.2,8.6);\node at (7.1,8.0) {例5：配方后直接看到两个固定点，交点在左侧};
\begin{scope}[shift={(5.5,2.2)},x=0.82cm,y=0.82cm]\draw[->,gray] (-4.6,0)--(2.6,0) node[right]{$x$};\draw[->,gray] (0,-.35)--(0,2.8) node[left]{$y$};\draw[orange!85!black,line width=1.6pt] (2,2)--(-4,0);\draw[violet,thick] (-1,1)--(-4,0);\draw[teal!80!black,thick] (2,2)--(-1,1);\fill[teal!80!black] (2,2) circle(2.2pt) node[above=6pt]{$A(2,2)$};\fill[teal!80!black] (-1,1) circle(2.2pt) node[above left=4pt]{$B(-1,1)$};\fill[blue] (-4,0) circle(2.2pt) node[below=4pt]{$P(-4,0)$};\end{scope}\node[anchor=west,align=left] at (8.8,4.3) {$A(2,2),\ B(-1,1)$\\[11pt]$P_0(-4,0)$\\[12pt]$d=3,\ D=\sqrt{10}$\\[12pt]$PA-PB=\sqrt{10}$\\[12pt]值域 $(-3,\sqrt{10}]$};\node at (7.1,7.3) {$\sqrt{x^2-4x+8}-\sqrt{x^2+2x+2}=\sqrt{(x-2)^2+4}-\sqrt{(x+1)^2+1}$};
\node[text width=12.8cm,align=center] at (7.1,.65) {图中横纵同尺度；点顺序为 $P_0,B,A$，所以距离差为正的 $AB$。};
\end{tikzpicture}\end{center}



	求
	\[
		f(x)=\sqrt{x^2-4x+8}-\sqrt{x^2+2x+2}
	\]
	的值域。

	\textbf{第 1 步：先配方。}
	\[
		x^2-4x+8=(x-2)^2+2^2,\qquad
		x^2+2x+2=(x+1)^2+1^2.
	\]
	所以 $m=2,n=-1,a=2,b=1$，属于 $a>b$。

	\textbf{第 2 步：算横向相隔多远，以及两点本身相隔多远。}
	\[
		d=|2-(-1)|=3,\quad D=\sqrt{3^2+1^2}=\sqrt{10}.
	\]
	\textbf{第 3 步：完整值域。}
	\[
		\boxed{f(x)\in(-3,\sqrt{10}]}.
	\]
	最大值在
	\[
		x_0=\frac{1\cdot2-2\cdot(-1)}{1-2}=-4
	\]
	处取得，代入得 $f(-4)=2\sqrt{10}-\sqrt{10}=\sqrt{10}$。
\end{examplebox}
\clearpage
\begin{examplebox}
	\textbf{\textcolor{CExample}{例 6：题目限制 $x>4$，原来的最大值也取不到}}
\begin{center}\begin{tikzpicture}[font=\small,>=stealth,line join=round]
\draw[gray!40,rounded corners=3pt] (0,0) rectangle (14.2,8.5);\node at (7.1,7.9) {例6：只差是否允许 x=4，值域右端开闭随之改变};
\begin{scope}[shift={(-0.2,2.3)},x=.48cm,y=17cm]\draw[->,gray] (3.4,0)--(12,0) node[right]{$x$};\draw[->,gray] (4,-.015)--(4,.28);\draw[gray,dashed] (3.5,0)--(11.5,0);\draw[blue,thick,domain=4.001:11.5,samples=130] plot(\x,{sqrt(\x*\x+4)-sqrt((\x-2)^2+1)-2});\draw[blue,fill=white] (4,{sqrt(5)-2}) circle(2.5pt);\node[below=4pt] at (4,0){$4$};\node[left=4pt] at (4,0){$2$};\node[blue,above right=8pt] at (4,{sqrt(5)-2}) {$\sqrt5$};\end{scope}\node at (3.1999999999999997,1.35) {$x>4$：$(2,\sqrt5)$};\begin{scope}[shift={(6.8,2.3)},x=.48cm,y=17cm]\draw[->,gray] (3.4,0)--(12,0) node[right]{$x$};\draw[->,gray] (4,-.015)--(4,.28);\draw[gray,dashed] (3.5,0)--(11.5,0);\draw[blue,thick,domain=4.001:11.5,samples=130] plot(\x,{sqrt(\x*\x+4)-sqrt((\x-2)^2+1)-2});\draw[blue,fill=blue] (4,{sqrt(5)-2}) circle(2.5pt);\node[below=4pt] at (4,0){$4$};\node[left=4pt] at (4,0){$2$};\node[blue,above right=8pt] at (4,{sqrt(5)-2}) {$\sqrt5$};\end{scope}\node at (10.2,1.35) {$x\ge4$：$(2,\sqrt5]$};
\node[text width=12.8cm,align=center] at (7.1,.65) {两图同尺度，纵轴从 $2$ 起显示实际函数值；下方虚线 $y=2$ 永远取不到。};
\end{tikzpicture}\end{center}



	仍取例 1 的函数：
	\[
		f(x)=\sqrt{x^2+4}-\sqrt{(x-2)^2+1},\qquad x>4.
	\]
	在整个实数范围内，它的最大值是 $\sqrt5$，且\textbf{只有 $x=4$ 时达到}。但现在题目明确说 $x>4$，因此 $x=4$ 不允许。

	导数告诉我们：$x>4$ 时函数不断\textbf{减小}；当 $x\to+\infty$ 时函数越来越接近 $2$，却不能等于 $2$。

	于是这道限制后的题\textbf{既没有最大值，也没有最小值}：
	\[
		\boxed{f(x)\in(2,\sqrt5)}.
	\]
	如果改成 $x\geq4$，那么 $x=4$ 就合法，右端立刻变成闭的：
	\[
		\boxed{f(x)\in(2,\sqrt5]}.
	\]
	\textbf{考试最容易丢分的位置就在这个括号上。}
\end{examplebox}
\clearpage
\begin{examplebox}
	\textbf{\textcolor{CExample}{例 7：闭区间上的值域，先找转折点再比较端点}}
\begin{center}\begin{tikzpicture}[font=\small,>=stealth,line join=round]
\draw[gray!40,rounded corners=3pt] (0,0) rectangle (14.2,8.8);\node at (7.1,8.200000000000001) {例7：只保留闭区间上的蓝线，两个端点都能取到};
\begin{scope}[shift={(3.2,3.0)},x=1cm,y=1.5cm]\draw[->,gray] (-1,0)--(5.4,0) node[right]{$x$};\draw[->,gray] (0,-.7)--(0,2.9) node[left]{$f$};\draw[gray,dashed,domain=-.7:5,samples=120] plot(\x,{sqrt(\x*\x+4)-sqrt((\x-2)^2+1)});\draw[blue,very thick,domain=0:4,samples=140] plot(\x,{sqrt(\x*\x+4)-sqrt((\x-2)^2+1)});\fill[blue] (0,{2-sqrt(5)}) circle(2.3pt);\fill[blue] (4,{sqrt(5)}) circle(2.3pt);\node at (.2,-1.05) {$2-\sqrt5$};\node[above left=5pt] at (4,{sqrt(5)}) {$\sqrt5$};\draw[gray,dotted] (4,0)--(4,{sqrt(5)});\node[below] at (4,0){$4$};\end{scope}\node[anchor=west,align=left] at (9.2,4.5) {$0\le x\le4$\\[12pt]蓝色保留整段递增部分\\[12pt]两个端点都允许\\[12pt]值域 $[2-\sqrt5,\sqrt5]$};
\node[text width=12.8cm,align=center] at (7.1,.65) {灰虚线仅提示原函数的延伸，不属于本题允许的图像。};
\end{tikzpicture}\end{center}



	仍考虑
	\[
		f(x)=\sqrt{x^2+4}-\sqrt{(x-2)^2+1},\qquad 0\leq x\leq4.
	\]
	因为此函数在 $x<4$ 时一直增大，$x=4$ 处最大，所以
	\[
		\begin{aligned}
			f(0) & =2-\sqrt5, \\
			f(4) & =\sqrt5.
		\end{aligned}
	\]
	因此完整值域为
	\[
		\boxed{[2-\sqrt5,\sqrt5]}.
	\]
	这次虽然还是同一个函数，但不再有无限远处的边界 $-2$，因为 $x$ 被限制在一小段路上。\textbf{受限范围优先看：是否包含转折点、左右端点是否允许、端点函数值谁大谁小。}
\end{examplebox}
\clearpage
\begin{examplebox}
	\textbf{\textcolor{CExample}{例 8：不用求导，直接求距离差为零的位置}}
\begin{center}\begin{tikzpicture}[font=\small,>=stealth,line join=round]
\draw[gray!40,rounded corners=3pt] (0,0) rectangle (14.2,8.7);\node at (7.1,8.1) {例8：根号差为零，就是到两定点一样远};
\begin{scope}[shift={(2.5,2.2)},x=1.5cm,y=1.5cm]\draw[->,gray] (-0.6,0)--(2.6,0) node[right]{$x$};\draw[->,gray] (0,-.35)--(0,2.8) node[left]{$y$};\draw[orange!85!black,line width=1.6pt] (0,2)--(0.25,0);\draw[violet,thick] (2,1)--(0.25,0);\draw[teal!80!black,thick] (0,2)--(2,1);\fill[teal!80!black] (0,2) circle(2.2pt) node[above left=4pt]{$A$};\fill[teal!80!black] (2,1) circle(2.2pt) node[above right=4pt]{$B$};\fill[blue] (0.25,0) circle(2.2pt) node[below=4pt]{$P$};\end{scope}\begin{scope}[shift={(2.5,2.2)},x=1.5cm,y=1.5cm]\draw[gray,dashed] (.1,-.3)--(1.55,2.6);\fill[gray] (1,1.5) circle(1.8pt);\node[above left=7pt] at (1,1.5){$M$};\draw[gray] (1.134,1.433)--(1.201,1.567)--(1.067,1.634);\end{scope}\node[anchor=west,align=left] at (8.5,4.4) {$PA=PB$\\[12pt]$P$ 在 $AB$ 的垂直平分线上\\[12pt]与 $x$ 轴交于 $x=1/4$\\[12pt]有理化后分子为 $4x-1$\\[10pt]所以零点仅为 $1/4$。};
\node[text width=12.8cm,align=center] at (7.1,.65) {橙、紫两段等长；灰点 $M$ 为 $AB$ 中点，灰虚线为垂直平分线。};
\end{tikzpicture}\end{center}



	仍取例 1：
	\[
		\sqrt{x^2+4}-\sqrt{(x-2)^2+1}=0.
	\]
	等价于两个根号相等，所以两边平方：
	\[
		x^2+4=(x-2)^2+1=x^2-4x+5.
	\]
	$x^2$ 抵消，只剩 $4=5-4x$，故
	\[
		\boxed{x=\frac14}.
	\]
	这里根号都不小于 $0$，所以从“两根号相等”变成“根号内相等”完全合法，不会凭空多出其他解。

	\textbf{更快的形式：}因为
	\[
		f(x)=\frac{x^2+4-\bigl((x-2)^2+1\bigr)}{
			\sqrt{x^2+4}+\sqrt{(x-2)^2+1}}
		=\frac{4x-1}{\sqrt{x^2+4}+\sqrt{(x-2)^2+1}},
	\]
	分母恒为正，所以只看分子 $4x-1$ 的正负就知道函数值的正负。
\end{examplebox}
